\section{实体识别}
\label{sec:nlp-extraction-entities}

实体识别从原文中找出指称对象的片段并给出类型，例如通知中的“海棠楼201”是地点，“海棠大学计算机学院”是组织。它的输出是本章其余任务的输入：实体链接为提及确定知识库身份，关系抽取在提及之间判断关系，事件抽取把提及填入论元角色。本节先说明实体结构有哪些形态、各路方法能表达哪些形态，再按结构与数据限制依次讨论序列标注、跨度、超图与词对网格、生成、少样本与提示等路线。

\subsection{实体结构的挑战与方法谱系}
\label{sec:nlp-extraction-mentions}

“海棠楼201”涉及三种不同的产出。它在原文中占据一段字符，这是\term{实体提及}（entity mention）；补充句中的“该教室”可能与它指向同一对象，这是文档内共指（第\ref{sec:nlp-analysis-span-coreference}小节）；把它对应到校园场所库中的某个编号，才是实体链接（第\ref{sec:nlp-extraction-linking}节）。图~\ref{fig:nlp-extraction-identity}用一段补充通知和一个小型知识库画出三者的关系。本节只处理第一种产出。

\begin{figure*}[htbp]
  \centering
  % Adapted from academic-drawing's TikZ template; project fonts and palette prevail.
% Final canvas: 167 mm wide. Ports: text east -> spans west; spans east
% -> cluster west; cluster east -> KB west. Side branches use a routing lane.
\begin{tikzpicture}[
  x=1mm,y=1mm,
  font=\figurefont\fontsize{9}{11}\selectfont,
  box/.style={rectangle,align=left,inner sep=2mm,line width=0.85pt,
    draw=BookRule,fill=white},
  source/.style={box,text width=31mm,draw=FigureInput,fill=FigureInput!8},
  span/.style={box,text width=27mm,draw=FigureInput,fill=FigureInput!8},
  cluster/.style={box,text width=23mm,draw=FigureModel,fill=FigureModel!8},
  target/.style={box,text width=33mm,draw=FigureOutput,fill=FigureOutput!8},
  locate/.style={draw=FigureInput,line width=0.85pt,rounded corners=1.2mm},
  coref/.style={draw=FigureModel,line width=0.85pt,dashed,rounded corners=1.2mm},
  link/.style={draw=FigureOutput,line width=1.1pt,-{Latex[length=2mm]}},
  head/.style={font=\figurefont\bfseries\fontsize{9}{11}\selectfont,
    text=BookInk,align=center}
]
  \node[head] at (18,37) {补充原文};
  \node[head] at (62,37) {原文跨度};
  \node[head] at (104,37) {文档共指簇};
  \node[head] at (146,37) {教学知识库};
  \node[source] (text) at (18,0) {
    地点为海棠楼201。\\
    该教室旁设临时服务点。\\[3pt]
    字符从0开始计数
  };
  \node[span] (m1) at (62,23) {海棠楼201\\$[3,9)$};
  \node[span] (m2) at (62,0) {该教室\\$[10,13)$};
  \node[span] (m3) at (62,-23) {临时服务点\\$[15,20)$};
  \node[cluster] (c1) at (104,11.5) {同一教室\\$\{m_1,m_2\}$};
  \node[cluster] (c2) at (104,-23) {服务点\\$\{m_3\}$};
  \node[target] (kb1) at (146,11.5) {ROOM-201\\海棠楼201};
  \node[target] (kb2) at (146,-23) {NIL\\无可确认库中项};
  \draw[locate] (text.east) -- (42,0);
  \draw[locate] (42,23) -- (42,-23);
  \draw[locate] (42,23) -- (m1.west);
  \draw[locate] (42,0) -- (m2.west);
  \draw[locate] (42,-23) -- (m3.west);
  \draw[coref] (m1.east) -- (84,23) -- (84,11.5) -- (c1.west);
  \draw[coref] (m2.east) -- (84,0) -- (84,11.5);
  \draw[coref] (m3.east) -- (c2.west);
  \draw[link] (c1.east) -- (kb1.west);
  \draw[link] (c2.east) -- (kb2.west);
  \draw[locate] (4,-38) -- (14,-38);
  \node[anchor=west] at (16,-38) {定位};
  \draw[coref] (49,-38) -- (59,-38);
  \node[anchor=west] at (61,-38) {共指归组};
  \draw[link] (111,-38) -- (121,-38);
  \node[anchor=west] at (123,-38) {身份选择};
\end{tikzpicture}

  \caption{实体提及、共指与身份链接。教学补充材料为“地点为海棠楼201。该教室旁设临时服务点”。原文到跨度的实线表示定位，跨度到共指簇的虚线表示同指归组，簇到知识库的箭头表示身份选择。临时服务点的NIL表示当前知识库没有可确认的对应项，不表示该对象不存在。}
  \label{fig:nlp-extraction-identity}
\end{figure*}

\term{命名实体识别}（named entity recognition, NER），本书简称实体识别，要同时回答“哪里是提及”和“属于什么类型”。沿用第\ref{chap:nlp-introduction}章的位置约定，原文字符区间写成从零开始的半开区间$[s,e)$；模型内部的词元从1开始编号，词元跨度$(i,j)$包含两端。若词元$i$对应原文$[a_i,b_i)$，连续跨度$(i,j)$恢复为$[a_i,b_j)$。子词切分和全半角转换都必须保留这张偏移表，否则模型给出的索引不能落回原文。

实体有四种结构形态，它们决定了输出空间。\emph{平面}实体互不重叠，一句话的实体集合可以写成一组不相交的区间。\emph{嵌套}实体允许一个区间包含另一个：“海棠大学计算机学院”是组织，其中“海棠大学”也是组织，这是\term{嵌套提及}（nested mention）。\term{不连续提及}（discontinuous mention）由几个不相邻的片段共同表达一个对象：在“左和右肺均清晰”中，若标注规范把省略了中心词的“左……肺”记为一个身体部位，它就由“左”和“肺”两段组成，输出应是有序片段列表，而不是覆盖中间文字的大区间。第四种困难与结构形式无关，而与输入单位有关：中文没有空格，字符模型不受分词错误影响，却丢失了“海棠楼”是一个词这类信息；词模型有词信息，但分词一旦切错，实体边界就无法恢复。

表~\ref{tab:nlp-extraction-ner-lineage}把本节讨论的方法按能表达的结构排开。从上到下，输出表示从逐位置标签变为区间、超图、词对网格，最后变为可线性化的序列，能表达的结构逐步变多；推断代价则从线性升到平方，生成方法还要承担串行解码与原文对齐。范式一列给出第\ref{sec:nlp-analysis-paradigms}节中的归类。

\begin{table*}[htbp]
  \centering
  \small
  \begin{tabularx}{\linewidth}{@{}>{\raggedright\arraybackslash}p{27mm}>{\raggedright\arraybackslash}p{15mm}*{3}{>{\centering\arraybackslash}p{10mm}}>{\centering\arraybackslash}p{12mm}>{\raggedright\arraybackslash}X@{}}
    \toprule
    方法 & 范式 & 平面 & 嵌套 & 不连续 & 词典信息 & 推断与代价\\
    \midrule
    BIO--CRF & 序列标注 & 能 & 否 & 否 & 不用 & 带类型掩码的Viterbi，$O(nL^2)$\\
    分层标注 & 序列标注 & 能 & 限层数 & 否 & 不用 & 逐层Viterbi，层数随嵌套深度增加\\
    Lattice-LSTM、FLAT & 序列标注 & 能 & 否 & 否 & 用 & Lattice-LSTM难并行；FLAT注意力为$O((n+m)^2)$\\
    双仿射跨度分类 & 跨度 & 能 & 能 & 否 & 不用 & $O(n^2)$候选，平面任务需冲突解码\\
    阅读理解式 & 跨度 & 能 & 能 & 否 & 不用 & 每个类型单独编码一次\\
    超图 & 图方法 & 能 & 能 & 扩展后有限 & 不用 & 超图inside递推，$O(n|T|)$；有结构歧义\\
    W2NER词对网格 & 图方法 & 能 & 能 & 能 & 不用 & $n\times n$网格分类，路径解码\\
    BARTNER、UIE & 约束生成 & 能 & 能 & 能 & 不用 & 串行解码，输出需解析并对齐原文\\
    \bottomrule
  \end{tabularx}
  \caption{实体识别的方法谱系。$n$为词元数，$L$为标签数，$|T|$为实体类型数，$m$为词典匹配到的词数。“词典信息”指是否显式接入中文词典词，不用词典的方法以字或子词为单位同样可以运行；“扩展后有限”指原始超图只表示连续提及，加入间隙状态的扩展可表示部分不连续提及。}
  \label{tab:nlp-extraction-ner-lineage}
\end{table*}

评价采用严格实体级F1：预测集合$\hat{\mathcal E}$与金标准集合$\mathcal E^\star$的元素是“片段列表加类型”，两者完全相同才算命中，$P=|\hat{\mathcal E}\cap\mathcal E^\star|/|\hat{\mathcal E}|$，$R=|\hat{\mathcal E}\cap\mathcal E^\star|/|\mathcal E^\star|$，$F_1=2PR/(P+R)$。把“海棠楼201”预测成“海棠楼”或预测成组织，都不计真阳性。细粒度和少样本设定另有基准：Few-NERD标注了8个粗类与66个细类，并按类型划分训练与测试，用来检验模型能否在少量样本下识别训练时没见过的细类\scite{ding2021fewnerd}。

\subsection{基于序列标注的实体识别}
\label{sec:nlp-extraction-ner-tagging}

序列标注是实体识别最早成熟的路线。输出表示、合法转移掩码、CRF训练与Viterbi解码已在第\ref{sec:nlp-analysis-paradigm-tagging}小节写出，BiLSTM-CRF的逐层运算见第\ref{sec:nlp-analysis-crf-chunk}小节；本小节只写实体任务特有的标签集、掩码与中文词汇增强。这条路线训练和解码都是线性时间，但一层标签不能表示嵌套与不连续，中文字符模型又缺少词边界信息；词汇增强针对后一问题。

\subsubsection{BIO--CRF与BiLSTM-CRF实体标注}
\label{sec:nlp-extraction-bio}

实体标注与分词、组块的差别在于标签带类型。设实体类型集合为$T$，BIO标签集为$\{\mathrm B\text{-}X,\mathrm I\text{-}X:X\in T\}\cup\{\mathrm O\}$，共$L=2|T|+1$个；BIOES再加E-$X$与S-$X$，共$4|T|+1$个。类型使掩码多了一条约束：延续标签必须与它所接的开头同类型。BIO规范下的转移掩码与起始掩码为
\begin{equation}
  \bm M_{a,b}=
  \begin{cases}
    -\infty, & b=\mathrm I\text{-}X\ \text{且}\ a\notin\{\mathrm B\text{-}X,\mathrm I\text{-}X\},\\
    0, & \text{其他},
  \end{cases}
  \qquad
  m^{\mathrm{start}}_{\mathrm I\text{-}X}=-\infty,\quad X\in T,
  \label{eq:nlp-extraction-crf}
\end{equation}
其余起始项与全部结束项为0。把它代入式~\eqref{eq:nlp-analysis-legal-mask}，整句得分、CRF负对数似然和Viterbi解码不需要任何改动：O之后不能直接接I-ORG，B-LOC之后也不能接I-ORG，这两类序列在配分函数和$\argmax$中都自动消失。掩码与标签集一样是规范的一部分，训练和解码必须使用同一份。

发射分数的来源经历了三代：特征CRF使用字形、词典与上下文窗口特征；BiLSTM-CRF用词嵌入、字符级表示与双向LSTM取代人工特征\scite{nlp-analysis-lample2016,v9a-ma2016}；预训练编码器再以上下文表示取代前两者，按首子词对齐收集词级表示（第\ref{sec:nlp-classical-bert-heads}小节）。三代模型的CRF层完全相同，改变的只是$\bm h_i$从哪里来。

以教学切分“海棠楼／201／开放”为例，正确标签是B-LOC、I-LOC、O，恢复“海棠楼201”$[0,6)$。若某位置的发射分数使逐位置取最大得到O、I-LOC、O，第二个标签不合法；掩码下的Viterbi只能在B-LOC、I-LOC、O与B-LOC、O、O等合法序列中选择。B-ORG、I-ORG、O边界正确而类型错误，B-LOC、O、O类型正确却漏了房间号，二者在严格F1中都不计命中。若输入只被切成一个词元“海棠楼201开放”，任何标签都表示不出正确边界，这时需要字符级输入或下一小节的词汇增强。

一层标签不能表示嵌套。Ju等的分层模型把多个平面标注层叠起来：第一层BiLSTM-CRF识别最内层实体，把每个识别出的实体内的词表示取平均、合并成一个单元，作为下一层的输入，再识别包含它的外层实体，直到某层不再输出实体为止\scite{ju2018layered}。它保留了线性链解码，代价是层间误差会传递，嵌套深度受层数限制，同一跨度属于多个类型的情形也难以表示。

\subsubsection{中文词汇增强：Lattice-LSTM与FLAT}
\label{sec:nlp-extraction-lattice}

中文实体识别有两种输入单位可选。词模型先分词再标注，若“海棠楼201开放”被切成“海棠楼／201开放”，实体的右边界落在第二个词内部，任何词级标签都恢复不出“海棠楼201”；字符模型没有这种传播错误，却看不到“海棠楼”是一个已知地名。第\ref{sec:nlp-analysis-pretrained-lattice}小节介绍了词格的思想：把词典匹配到的词作为额外单元接入字符序列，让语境决定采用哪个候选词。这里写出两种实现的完整计算。

第一种实现是Zhang与Yang的Lattice-LSTM，它在字符LSTM上加入词单元\scite{nlp-analysis-zhang2018lattice}。设句子有$n$个字，第$j$个字的嵌入为$\bm x^{\mathrm c}_j\in\mathbb R^{d_{\mathrm c}}$，隐藏维度为$d_{\mathrm h}$；词典$D$中从第$b$个字到第$e$个字的匹配词记为$w_{b,e}$，嵌入为$\bm x^{\mathrm w}_{b,e}\in\mathbb R^{d_{\mathrm w}}$；以第$j$个字结尾的匹配词起点集合为$B_j=\{b:w_{b,j}\in D\}$。字符单元先按普通LSTM算出各门与候选更新：
\begin{equation}
  \begin{bmatrix}\bm i^{\mathrm c}_j\\ \bm o^{\mathrm c}_j\\ \bm f^{\mathrm c}_j\\ \tilde{\bm c}^{\mathrm c}_j\end{bmatrix}
  =\begin{bmatrix}\sigma\\ \sigma\\ \sigma\\ \tanh\end{bmatrix}
  \left(\bm W^{\mathrm c}\begin{bmatrix}\bm x^{\mathrm c}_j\\ \bm h^{\mathrm c}_{j-1}\end{bmatrix}+\bm b^{\mathrm c}\right),
  \qquad \bm W^{\mathrm c}\in\mathbb R^{4d_{\mathrm h}\times(d_{\mathrm c}+d_{\mathrm h})}.
  \label{eq:nlp-extraction-lattice-char}
\end{equation}
每个匹配词$w_{b,e}$有自己的词单元。它以词嵌入和起点字的隐藏状态为输入，没有输出门，因为词单元不直接产生隐藏状态，只向词尾字输送记忆：
\begin{equation}
  \begin{bmatrix}\bm i^{\mathrm w}_{b,e}\\ \bm f^{\mathrm w}_{b,e}\\ \tilde{\bm c}^{\mathrm w}_{b,e}\end{bmatrix}
  =\begin{bmatrix}\sigma\\ \sigma\\ \tanh\end{bmatrix}
  \left(\bm W^{\mathrm w}\begin{bmatrix}\bm x^{\mathrm w}_{b,e}\\ \bm h^{\mathrm c}_{b}\end{bmatrix}+\bm b^{\mathrm w}\right),
  \qquad
  \bm c^{\mathrm w}_{b,e}=\bm f^{\mathrm w}_{b,e}\odot\bm c^{\mathrm c}_{b}+\bm i^{\mathrm w}_{b,e}\odot\tilde{\bm c}^{\mathrm w}_{b,e},
  \label{eq:nlp-extraction-lattice-word}
\end{equation}
其中$\bm W^{\mathrm w}\in\mathbb R^{3d_{\mathrm h}\times(d_{\mathrm w}+d_{\mathrm h})}$。在词尾字$j$处，每个词单元再得到一个额外输入门$\bm i^{\mathrm c}_{b,j}=\sigma(\bm W^{\mathrm l}[\bm x^{\mathrm c}_j;\bm c^{\mathrm w}_{b,j}]+\bm b^{\mathrm l})$，它衡量这个词对当前字有多大贡献。字符单元的记忆由词单元与自身候选按门控加权合并：
\begin{equation}
  \bm c^{\mathrm c}_j=\sum_{b\in B_j}\bm\alpha^{\mathrm c}_{b,j}\odot\bm c^{\mathrm w}_{b,j}+\bm\alpha^{\mathrm c}_j\odot\tilde{\bm c}^{\mathrm c}_j,
  \qquad
  \bm\alpha^{\mathrm c}_{b,j}=\frac{\exp(\bm i^{\mathrm c}_{b,j})}{\exp(\bm i^{\mathrm c}_j)+\sum_{b'\in B_j}\exp(\bm i^{\mathrm c}_{b',j})},
  \label{eq:nlp-extraction-lattice-merge}
\end{equation}
$\bm\alpha^{\mathrm c}_j$的分子换成$\exp(\bm i^{\mathrm c}_j)$，指数与除法都逐维进行，因而每一维上全部权重之和为1。最后$\bm h^{\mathrm c}_j=\bm o^{\mathrm c}_j\odot\tanh(\bm c^{\mathrm c}_j)$。若$B_j$为空，字符单元退回普通LSTM的$\bm c^{\mathrm c}_j=\bm f^{\mathrm c}_j\odot\bm c^{\mathrm c}_{j-1}+\bm i^{\mathrm c}_j\odot\tilde{\bm c}^{\mathrm c}_j$。注意有匹配词时式中不再出现$\bm c^{\mathrm c}_{j-1}$：前文信息改由词单元从起点$b$带过来。$\bm h^{\mathrm c}_j$之上接发射层与带类型掩码的CRF，整网用式~\eqref{eq:nlp-analysis-legal-mask}的负对数似然端到端训练。

以“海棠楼201”与教学词典$\{$海棠，海棠楼$\}$为例，匹配词为$w_{1,2}$与$w_{1,3}$，两个词单元都从第1个字的状态出发，分别在“棠”和“楼”处并入。看“楼”处的某一维，设$i^{\mathrm c}_3=0.3$、$i^{\mathrm c}_{1,3}=0.9$，则$\alpha^{\mathrm c}_{1,3}=e^{0.9}/(e^{0.3}+e^{0.9})\approx0.6457$，$\alpha^{\mathrm c}_3\approx0.3543$；若词单元在这一维的记忆为0.8、字符候选为$-0.5$，合并结果为$0.6457\times0.8-0.3543\times0.5\approx0.3394$，主要由“海棠楼”的词信息决定。这个算例也显示出式~\eqref{eq:nlp-extraction-lattice-merge}的一个性质：门值经sigmoid后落在$(0,1)$，再取指数归一化，两个候选的权重只能在$1/(1+e)\approx0.27$与$e/(1+e)\approx0.73$之间，任何一方都不会被完全关掉。

Lattice-LSTM的代价来自结构本身。每个字的输入数随匹配词数变化，词单元又要等起点字算完才能计算，整句只能按字串行推进，难以在批内对齐并行；词典的覆盖与质量也直接决定可用的词信息。

第二种实现针对这一代价。Li等的FLAT（flat-lattice Transformer）保留词格的信息，换掉串行结构\scite{li2020flat}。它把字与匹配词放进同一个序列：$n$个字在前，$m$个匹配词在后，每个单元$u$记录它在原句中的头位置$\mathrm{head}[u]$与尾位置$\mathrm{tail}[u]$；字的头尾相同，词的头尾是首字与末字的位置。“海棠楼201”展平后共8个单元：6个字的头尾为$(1,1),\ldots,(6,6)$，“海棠”为$(1,2)$，“海棠楼”为$(1,3)$。序列中的先后次序不再表示原句位置，位置信息全部由头尾坐标给出。

两个单元之间的关系由四个相对距离刻画：
\begin{equation}
  \begin{aligned}
    d^{(\mathrm{hh})}_{u,v}&=\mathrm{head}[u]-\mathrm{head}[v],&\quad
    d^{(\mathrm{ht})}_{u,v}&=\mathrm{head}[u]-\mathrm{tail}[v],\\
    d^{(\mathrm{th})}_{u,v}&=\mathrm{tail}[u]-\mathrm{head}[v],&\quad
    d^{(\mathrm{tt})}_{u,v}&=\mathrm{tail}[u]-\mathrm{tail}[v].
  \end{aligned}
  \label{eq:nlp-extraction-flat-distance}
\end{equation}
四个数合起来能区分包含、相交与相邻。“海棠楼”与“棠”的四个距离为$(-1,-1,1,1)$，头在其左、尾在其右，表示包含；“海棠”与“楼”为$(-2,-2,-1,-1)$，尾距为$-1$，表示前者紧接在后者之前。每个距离$d$先映射为正弦位置向量$\bm p_d$（$p_d^{(2k)}=\sin(d/10000^{2k/d_{\mathrm{model}}})$，$p_d^{(2k+1)}$取余弦），再融合为相对位置表示$\bm R_{u,v}=\operatorname{ReLU}(\bm W_{\mathrm r}[\bm p_{d^{(\mathrm{hh})}};\bm p_{d^{(\mathrm{th})}};\bm p_{d^{(\mathrm{ht})}};\bm p_{d^{(\mathrm{tt})}}])$。注意力分数采用Transformer-XL式的相对位置写法，设$\bm e_u$为单元嵌入，$\bm q_u=\bm W_{\mathrm q}\bm e_u$，$\bm k_v=\bm W_{\mathrm k}\bm e_v$：
\begin{equation}
  A_{u,v}=\bm q_u^\top\bm k_v+\bm q_u^\top\bm W_{\mathrm R}\bm R_{u,v}+\bm a^\top\bm k_v+\bm c^\top\bm W_{\mathrm R}\bm R_{u,v},
  \label{eq:nlp-extraction-flat-attention}
\end{equation}
四项依次是内容与内容、内容与相对位置、全局内容偏置与全局位置偏置，$\bm a,\bm c$为可学习向量。分数在全部$n+m$个单元上做softmax，其余与标准自注意力相同；自注意力的一般计算可参考\BookChapterRef{chap:transformer}{卷五的Transformer章}。最后只取前$n$个字单元的输出接CRF，词单元只在注意力中提供信息。

与Lattice-LSTM相比，FLAT把“词在哪里、和哪个字是什么关系”从网络拓扑移到了位置编码里。拓扑一旦不再随句子变化，整句可以一次并行计算；每个字也能直接注意到所有匹配词，而不仅是以自己结尾的词。代价是注意力的计算与内存为$O((n+m)^2)$，匹配词多时序列明显变长；两者都依赖词典，词典里没有的新名称得不到词信息。预训练编码器普及后，字级BERT已经从大规模语料中学到不少词边界知识，词汇增强在领域词典质量高、标注语料小的场景中更有价值。

\subsection{基于跨度的实体识别}
\label{sec:nlp-extraction-ner-span}

为表示嵌套实体，跨度路线直接对起止对分类，跨度表示与冲突解码见第\ref{sec:nlp-analysis-paradigm-span}小节。它允许候选相互重叠，代价是$O(n^2)$个候选和平面任务中额外的冲突选择。

\subsubsection{跨度分类与双仿射起止评分}
\label{sec:nlp-extraction-biaffine}

候选集合为$C=\{(i,j):1\leq i\leq j\leq n\}$，每个候选取一个实体类型或“非实体”$\varnothing$。最直接的打分器把式~\eqref{eq:nlp-analysis-span-representation}的跨度表示送入前馈网络。Yu等改用依存分析中的双仿射打分\scite{yu2020ner}：编码器（BERT、词与字符表示加BiLSTM）给出$\bm h_i$后，两个前馈网络分别把它投影为起点表示$\bm p_i=f_{\mathrm s}(\bm h_i)$与终点表示$\bm q_j=f_{\mathrm e}(\bm h_j)\in\mathbb R^{k}$，类型$t$的分数为
\begin{equation}
  z_{i,j,t}=
  \bm p_i^\top\bm U_t\bm q_j+
  \bm w_t^\top[\bm p_i;\bm q_j]+b_t,
  \qquad
  p(t\mid i,j,x)=
  \frac{\exp z_{i,j,t}}{\sum_{t'\in T\cup\{\varnothing\}}\exp z_{i,j,t'}},
  \label{eq:nlp-extraction-biaffine}
\end{equation}
其中$\bm U_t\in\mathbb R^{k\times k}$、$\bm w_t\in\mathbb R^{2k}$、$b_t\in\mathbb R$。双线性项刻画两端是否配对，仿射项保留单独起点和终点的倾向。把全部类型叠在一起，$\bm U$是$k\times(|T|+1)\times k$的张量，一次运算就得到$n\times n\times(|T|+1)$的分数张量，只取上三角$i\leq j$。训练时金标准跨度取其类型，其余候选取$\varnothing$，最小化$L_{\mathrm{span}}=-\sum_{(i,j)\in C}\log p(t^{*}_{i,j}\mid i,j,x)$；$\varnothing$远多于正例，可降采样或降权，但评价必须在完整候选集上进行。

\begin{example}[嵌套候选怎样保留]
\label{ex:nlp-extraction-nested}
教学词元序列为“海棠／大学／计算机／学院／发布／通知”，只用一个实体类别ORG，非实体分数均为0。候选及其ORG分数为$(1,4):4.0$、$(1,2):3.0$、$(2,5):2.0$、$(3,4):1.5$。取$(1,4)$的$\bm p_1=(1,0)^\top$、$\bm q_4=(2,1)^\top$、$\bm U_{\mathrm{ORG}}=\bm I_2$、$\bm w_{\mathrm{ORG}}=(0,0,1,0)^\top$、$b_{\mathrm{ORG}}=0$，式~\eqref{eq:nlp-extraction-biaffine}给出$2+2=4$，实体概率为$\exp(4)/(1+\exp(4))\approx0.9820$，该候选损失约为0.0181。

按分数降序扫描，先保留$(1,4)$。嵌套规范保留其内部的$(1,2)$和$(3,4)$，排除与它部分交叉的$(2,5)$，因为$1<2\leq4<5$；平面规范只保留$(1,4)$。对应字符区间为$[0,9)$、$[0,4)$与$[4,9)$。贪心扫描不保证总分最大，平面规范下的精确解是第\ref{sec:nlp-analysis-paradigm-span}小节的加权区间调度。
\end{example}

与独立的起点、终点指针相比，双仿射给每个起止组合单独评分。若只预测起点集合$\{1,3\}$和终点集合$\{2,4\}$，无法判断$(1,4)$是否成立；双仿射直接对$(1,4)$打分，消除了这种配对歧义。它仍有两点限制：候选为平方级，长文档常需设最大宽度；每个$(i,j)$仍是连续区间，表示不了不连续提及。

\subsubsection{查询式阅读理解NER}
\label{sec:nlp-extraction-mrc-ner}

类型标签“组织”本身携带的信息很少，模型只能从标注样本里猜它的边界。Li等把标注指南改写为查询，例如“找出学校、学院等组织名称”，与正文拼接后送入编码器，把实体识别改成抽取式阅读理解：查询指定找哪类对象，答案是原文中的跨度\scite{li2020mrcner}。每个类型提一次问，不同查询的答案可以重叠，于是不同类型之间的嵌套自然得到保留。

设类型$t$的查询为$c_t$，编码$[\mathrm{CLS}]\,c_t\,[\mathrm{SEP}]\,x\,[\mathrm{SEP}]$后取正文部分的表示$\bm H^{(t)}\in\mathbb R^{n\times d}$。一篇文本可能有多个同类实体，因此模型不是在全文中各选一个起点和终点，而是对每个位置独立做“是否起点”“是否终点”的二分类，两个$d\times2$的分类头给出逐行概率。起点集合与终点集合之间还要配对，匹配头为
\begin{equation}
  p_{i,j}^{(t)}=\sigma\!\left(
  \bm v^{\top}[\bm h^{(t)}_i;\bm h^{(t)}_j]\right),
  \qquad
  L=\alpha L_{\mathrm{s}}+\beta L_{\mathrm{e}}+\gamma L_{\mathrm{match}},
  \label{eq:nlp-extraction-mrc}
\end{equation}
其中$\bm v\in\mathbb R^{2d}$，三个损失分别为起点、终点与配对的二元交叉熵，权重在开发集上选择。没有该类型实体的文本提供全负的起止与配对标签。推理时在预测的起点与终点中枚举$i\leq j$的组合，经匹配头筛选后恢复字符偏移。组织查询应返回“海棠大学”和“海棠大学计算机学院”，地点查询在这句话上返回空集。

查询中的定义和例词让同一个编码器在不同类型间共享知识，适合类型多、需要说明含义的场景。代价是$|T|$个类型要编码$|T|$遍；同一类型内部的嵌套仍取决于匹配头能否正确配对，查询写法也会影响结果，比较时应固定查询模板。

\subsection{嵌套与不连续实体}
\label{sec:nlp-extraction-ner-nested}

不连续实体既不能用BIO表示，也不能用连续跨度表示。“左和右肺均清晰”中，“左……肺”与“右肺”共享中心词“肺”，前者的两段之间隔着“和右”。处理这类结构需要一种输出表示，既能表达“哪些词属于同一个实体”，又不要求这些词相邻。下面的超图与词对网格是两种做法。

\subsubsection{超图表示}
\label{sec:nlp-extraction-hypergraph}

一句话中可能的嵌套实体组合随长度指数增长，逐个枚举不可行。Lu与Roth提出\term{提及超图}（mention hypergraph），用多项式规模的节点与超边共同编码这些组合\scite{lu2015hypergraph}。超图中的一条\term{超边}（hyperedge）从一个父节点指向一组子节点，表示父节点的状态由这组子节点共同展开。下面按其思想给出一种简化写法。对每个位置$k$与类型$t$设五类节点：$A_k$表示“处理位置$k$及其后的所有提及”；$E_k$表示“从$k$开始的提及”；$T_k^t$表示“从$k$开始的类型$t$提及”；$I_k^t$表示“位置$k$在某个类型$t$提及内部”；终止节点$X$表示“到此结束”。超边规定节点可以怎样展开：$A_k\to(E_k,A_{k+1})$；$E_k\to(T_k^{1},\ldots,T_k^{|T|})$；$T_k^t\to X$表示没有类型$t$的提及从$k$开始，$T_k^t\to I_k^t$表示有；$I_k^t$有三种展开，$\to I_{k+1}^t$表示提及延续，$\to X$表示提及在$k$结束，$\to(I_{k+1}^t,X)$表示一个提及在$k$结束、另一个更长的同时延续。

一个输入的标注结构是从$A_1$出发、每个访问到的节点各选一条超边形成的子超图，提及就是从某个$T_k^t$经若干$I^t$节点到达$X$的路径。以“海棠／大学／计算机／学院”的三个ORG提及$(1,4)$、$(1,2)$、$(3,4)$为例：$T_1\to I_1$，$I_1\to I_2$，$I_2\to(I_3,X)$表示$(1,2)$在第2个词结束而$(1,4)$继续，$T_3\to I_3$，$I_3\to I_4$，$I_4\to X$，其余$T_k\to X$。从$T_1$出发有两条到$X$的路径，读出$(1,2)$与$(1,4)$；从$T_3$出发读出$(3,4)$。$I_3$被两个提及共享，这正是超图紧凑的来源。

模型是超图上的对数线性模型。每条超边$\varepsilon$有特征$\bm f(\varepsilon,x)$，结构$y$的得分是所选超边得分之和，条件概率按全部合法子超图归一化。配分函数用超图上的inside递推计算：
\begin{equation}
  \begin{aligned}
    Z(X)&=1,\qquad
    Z(v)=\sum_{\varepsilon\in\mathrm{out}(v)}\exp\bigl(\bm w^\top\bm f(\varepsilon,x)\bigr)\prod_{u\in\mathrm{child}(\varepsilon)}Z(u),\\
    \log p(y\mid x)&=\sum_{\varepsilon\in y}\bm w^\top\bm f(\varepsilon,x)-\log Z(A_1),
  \end{aligned}
  \label{eq:nlp-extraction-hypergraph-inside}
\end{equation}
$\mathrm{out}(v)$为从$v$出发的超边，$\mathrm{child}(\varepsilon)$为其子节点。节点按位置从右向左计算，每个位置有$O(|T|)$个节点和常数条超边，总代价为$O(n|T|)$，与句长成线性；把求和换成取最大即得解码。这与第\ref{sec:nlp-analysis-paradigm-graph}小节的区间动态规划同属一类，只是递推定义在超图而非区间表上。后续工作把超边特征换成LSTM表示，训练目标不变。

线性复杂度的代价是\emph{结构歧义}：一个子超图可能对应不止一个提及集合。设同类型提及为$(1,3)$与$(2,4)$，超图需要$T_1\to I_1$、$I_1\to I_2$、$T_2\to I_2$、$I_2\to I_3$、$I_3\to(I_4,X)$、$I_4\to X$，从中可读出$(1,3)$、$(1,4)$、$(2,3)$、$(2,4)$四条路径，与“恰有$(1,3)$和$(2,4)$”无法区分。解码时只能约定一种读法，例如输出全部路径，这在同类型交叉提及上会产生多余实体。原始超图只表示连续提及；Muis与Lu在超图中为提及内部的词加入“属于实体”与“间隙”两种状态，从而能表示部分不连续实体，但歧义也随之增多\scite{muis2016discontiguous}。

\subsubsection{W2NER词元关系图}
\label{sec:nlp-extraction-w2ner}

超图把结构编码在节点状态里，需要为不同结构精心设计状态。W2NER换一个角度：既然不连续实体的困难在于“哪些词属于同一实体”，就直接预测词与词之间的关系\scite{li2022w2ner}。它定义两类有向关系：NNW（next-neighboring-word）从实体中的一个词指向该实体中次序上的下一个词，这里的“下一个”指实体内部的次序，不要求原文相邻；THW-$t$（tail-head-word）从实体的尾词指回首词，并携带类型$t$。对$n$个词，输出是$n\times n$的关系网格$\bm G$，每格取$\{\mathrm{NONE},\mathrm{NNW}\}\cup\{\mathrm{THW}\text{-}t:t\in T\}$之一。由于实体内的词按原文顺序排列，NNW总在上三角$i<j$，THW总在下三角或对角线$i\geq j$，单词实体就是对角线上的THW。

“左／和／右／肺／均／清晰”中，“左……肺”由第1、4个词组成，“右肺”由第3、4个词组成。正标签为$G_{1,4}=G_{3,4}=\mathrm{NNW}$、$G_{4,1}=G_{4,3}=\mathrm{THW}\text{-}\mathrm{BODY}$，其余为NONE，如图~\ref{fig:nlp-extraction-w2ner-grid}。两个实体共享尾词“肺”，在网格中分别占用不同的格，不会互相覆盖；平面与嵌套实体同样只是网格中的若干NNW路径和THW格。

\begin{figure}[htbp]
  \centering
  % 第九卷 10.1.4.2 节：W2NER词对关系网格（教学句“左／和／右／肺／均／清晰”）。
% 行为i、列为j；上三角(j>i)放NNW，下三角(i>j)放THW-类型，对角线可放单词实体的THW。
% 正标签：G(1,4)=G(3,4)=NNW，G(4,1)=G(4,3)=THW-BODY，其余为NONE。
\begin{tikzpicture}[
  x=1mm,y=1mm,figure picture,
  cellbox/.style={draw=FigureGrid,line width=.4pt,minimum size=8mm,inner sep=0pt},
  upper/.style={cellbox,fill=white},
  lower/.style={cellbox,fill=FigurePaper},
  nnw/.style={cellbox,fill=FigureModelFill,draw=FigureModel,line width=.7pt,
    font=\figurefont\scriptsize},
  thw/.style={cellbox,fill=FigureInputFill,draw=FigureInput,line width=.7pt,
    font=\figurefont\tiny,align=center},
  lab/.style={font=\figurefont\footnotesize,text=FigureInk},
  note/.style={font=\figurefont\scriptsize,text=FigureMuted,anchor=west,align=left}
]
  % 网格底色
  \foreach \i in {1,...,6}{
    \foreach \j in {1,...,6}{
      \pgfmathtruncatemacro{\up}{\j>\i ? 1 : 0}
      \pgfmathtruncatemacro{\dn}{\i>\j ? 1 : 0}
      \ifnum\up=1
        \node[upper] at ({(\j-1)*8},{-(\i-1)*8}) {};
      \else\ifnum\dn=1
        \node[lower] at ({(\j-1)*8},{-(\i-1)*8}) {};
      \else
        \node[cellbox,fill=white] at ({(\j-1)*8},{-(\i-1)*8}) {};
      \fi\fi
    }
  }
  % 行列标签
  \foreach \k/\w in {1/左,2/和,3/右,4/肺,5/均,6/清晰}{
    \node[lab] at ({(\k-1)*8},7) {\w};
    \node[font=\figurefont\scriptsize,text=FigureMuted] at ({(\k-1)*8},11.5) {$j{=}\k$};
    \node[lab,anchor=east] at (-6,{-(\k-1)*8}) {\w};
    \node[font=\figurefont\scriptsize,text=FigureMuted,anchor=east] at (-13,{-(\k-1)*8}) {$i{=}\k$};
  }
  % 正标签
  \node[nnw] at ({(4-1)*8},{-(1-1)*8}) {NNW};
  \node[nnw] at ({(4-1)*8},{-(3-1)*8}) {NNW};
  \node[thw] at ({(1-1)*8},{-(4-1)*8}) {THW\\BODY};
  \node[thw] at ({(3-1)*8},{-(4-1)*8}) {THW\\BODY};
  % 区域说明
  \node[note] at (50,0) {上三角（$j>i$）：NNW，\\实体内部的次序相邻词};
  \node[note] at (50,-14) {下三角（$i>j$）：THW-类型，\\由尾词指回首词};
  \node[note] at (50,-28) {解码：THW(4,1)与路径1$\to$4\\得“左…肺”$[0,1)\cup[3,4)$；};
  \node[note] at (50,-38) {THW(4,3)与路径3$\to$4\\得“右肺”$[2,4)$};
\end{tikzpicture}

  \caption{W2NER在“左和右肺均清晰”上的词对关系网格。第$i$行第$j$列表示从词$i$到词$j$的关系，未标注的格为NONE。上三角的NNW连接同一实体内次序相邻的词，下三角的THW-BODY由尾词指回首词并给出类型。两条NNW路径$1\to4$与$3\to4$共享尾词“肺”，分别由THW$(4,1)$与THW$(4,3)$确认。标注规范为教学设定。}
  \label{fig:nlp-extraction-w2ner-grid}
\end{figure}

网格中每一格需要一个词对表示，W2NER的逐层运算如下。
\begin{enumerate}
  \item 编码：BERT输出经首子词对齐与BiLSTM，得到$\bm H\in\mathbb R^{n\times d_{\mathrm h}}$。
  \item \term{条件层归一化}（conditional layer normalization, CLN）：用词$i$生成归一化的缩放与平移参数，作用在词$j$上，
  \[
    \bm V_{i,j}=\bm\gamma_{i}\odot\frac{\bm h_j-\mu_j}{\sigma_j}+\bm\lambda_{i},\qquad
    \bm\gamma_i=\bm W_{\gamma}\bm h_i+\bm b_{\gamma},\quad
    \bm\lambda_i=\bm W_{\lambda}\bm h_i+\bm b_{\lambda},
  \]
  $\mu_j,\sigma_j$是$\bm h_j$各维的均值与标准差。普通层归一化的缩放与平移是全局参数，这里由另一个词决定，于是$\bm V_{i,j}$表示“以$i$为条件看$j$”，得到$\bm V\in\mathbb R^{n\times n\times d_{\mathrm h}}$。
  \item 距离与区域嵌入：把$|i-j|$分桶后查表得距离嵌入$\bm E^{\mathrm d}_{i,j}$，再用区域嵌入$\bm E^{\mathrm r}_{i,j}$区分上三角、下三角与对角线；三者拼接后经多层感知机降维，得到$\bm C\in\mathbb R^{n\times n\times d_{\mathrm c}}$。区域嵌入让模型知道某格只可能是NNW还是THW。
  \item 多粒度空洞卷积：在网格上做空洞率$l\in\{1,2,3\}$的二维卷积，$\bm Q^{(l)}=\operatorname{GELU}(\operatorname{DConv}_l(\bm C))$，拼接得$\bm Q\in\mathbb R^{n\times n\times3d_{\mathrm c}}$。\term{空洞卷积}（dilated convolution）在卷积核元素之间留出间隔，不增加参数就扩大感受野；在这里它让一格能看到同行、同列附近的格，例如判断$G_{4,1}$时参考$G_{1,4}$是否为NNW。
  \item 联合预测：双仿射分支以编码器输出为输入，$\bm s_i=\operatorname{MLP}_{\mathrm s}(\bm h_i)$、$\bm o_j=\operatorname{MLP}_{\mathrm o}(\bm h_j)$，$\bm y'_{i,j}=\bm s_i^\top\bm U\bm o_j+\bm W[\bm s_i;\bm o_j]+\bm b$；多层感知机分支以网格表示为输入，$\bm y''_{i,j}=\operatorname{MLP}(\bm Q_{i,j})$。两者相加后逐格softmax，$p(\cdot\mid i,j)=\operatorname{softmax}(\bm y'_{i,j}+\bm y''_{i,j})\in\mathbb R^{|T|+2}$。
\end{enumerate}
训练对全部$n^2$个格做交叉熵，NONE格也参与。双仿射分支直接看两个词本身，网格分支看邻近格的关系模式，两者互补。

解码把网格还原为实体，步骤见算法~\ref{alg:nlp-extraction-w2ner-decode}。NNW边只从小下标指向大下标，构成有向无环图，从首词到尾词的路径可以用深度优先搜索枚举。

\Needspace{10\baselineskip}
\begin{algorithm}[W2NER网格解码]
  \label{alg:nlp-extraction-w2ner-decode}
  \begin{algorithmic}[1]
    \Require 预测网格$\hat{\bm G}\in\mathcal R^{n\times n}$，$\mathcal R=\{\mathrm{NONE},\mathrm{NNW}\}\cup\{\mathrm{THW}\text{-}t\}$；词到字符的偏移表
    \Ensure 实体集合$\hat{\mathcal E}$，每个实体为片段列表加类型
    \State 建有向图$\mathcal D$：对每个$i<j$且$\hat G_{i,j}=\mathrm{NNW}$，加边$i\to j$
    \State $\hat{\mathcal E}\gets\varnothing$
    \For{每个$i\geq j$且$\hat G_{i,j}=\mathrm{THW}\text{-}t$} \Comment{$i$为尾词，$j$为首词}
      \If{$i=j$}
        \State 加入单词实体$(\{j\},t)$
      \Else
        \For{$\mathcal D$中每条从$j$到$i$的路径$j=k_1\to k_2\to\cdots\to k_r=i$}
          \State 加入实体$(\{k_1,\ldots,k_r\},t)$
        \EndFor
      \EndIf
    \EndFor
    \State 把每个实体中下标连续的词合并为片段，按偏移表转为字符区间
  \end{algorithmic}
\end{algorithm}

在教学例子中，THW$(4,1)$要求从1到4的路径，图中只有边$1\to4$，得到$\{1,4\}$，两段字符区间为$[0,1)$与$[3,4)$；THW$(4,3)$得到路径$3\to4$，合并为一段$[2,4)$。“左和右肺”整体不会被当作一个实体，因为图中没有$1\to2$或$2\to3$的边。

W2NER用一套标签统一了三种结构，解码也只是图上找路径。它的弱点在两处。其一，网格大小为$n^2$，长文本的计算与内存压力明显，通常按句处理。其二，路径组合可能产生标注中不存在的实体：若两个同类型实体的NNW边交错，THW确认的首尾之间可能有多条路径，算法会把它们都输出。不同语料对省略形式是否算不连续实体的规定不一，比较结果前需要统一标注规范。

\subsection{基于生成的实体识别}
\label{sec:nlp-extraction-ner-generation}

前几小节的每种方法都为特定结构设计了专用标签和解码器：BIO管平面，跨度管嵌套，网格管不连续。生成路线把实体序列化为目标序列，一个序列到序列模型同时覆盖三种结构，并能直接利用预训练的编码器—解码器；线性化与约束解码的通用讨论见第\ref{sec:nlp-analysis-paradigm-generation}小节。

\subsubsection{BARTNER指针式生成}
\label{sec:nlp-extraction-bartner}

若让模型直接生成实体字符串，同一字符串在原文中出现多次时无法确定位置，模型还可能写出原文没有的字。Yan等的BARTNER让目标序列只含两种符号：指向输入位置的指针和类型标记\scite{yan2021bartner}。每个实体写成其各片段起止位置的指针序列，后接一个类型标记；多个实体按原文位置依次排列。“左和右肺均清晰”的目标为
\[
  \underbrace{1\ 1\ 4\ 4\ \langle\mathrm{BODY}\rangle}_{\text{左……肺}}\quad
  \underbrace{3\ 4\ \langle\mathrm{BODY}\rangle}_{\text{右肺}}\quad\langle\mathrm{eos}\rangle,
\]
第一个实体有两段$[1,1]$与$[4,4]$，故有四个指针。连续实体只有一对起止指针，嵌套实体就是指针范围相互包含的两个实体，三种结构用同一种序列表示。

解码器每步的输出空间是$n$个输入位置与$|C|$个类型标记（含结束符）的并集。设BART编码器输出$\bm H^{\mathrm e}\in\mathbb R^{n\times d}$，输入词元嵌入为$\bm E^{\mathrm e}\in\mathbb R^{n\times d}$，类型标记的嵌入为$\bm G\in\mathbb R^{|C|\times d}$。已生成的目标符号先换回词元：指针$y_u\leq n$换成输入的第$y_u$个词元，类型标记换成对应的类型词，这样解码器的输入仍是普通词元。第$u$步的解码器状态为$\bm g_u\in\mathbb R^{d}$，输出分布为
\begin{equation}
  \begin{aligned}
    \bar{\bm H}^{\mathrm e}&=\kappa\operatorname{MLP}(\bm H^{\mathrm e})+(1-\kappa)\bm E^{\mathrm e},\\
    p_\theta(y_u\mid\bm y_{<u},x)&=\operatorname{softmax}\bigl([\bar{\bm H}^{\mathrm e}\bm g_u;\ \bm G\bm g_u]\bigr)\in\mathbb R^{n+|C|},
  \end{aligned}
  \label{eq:nlp-extraction-bartner}
\end{equation}
前$n$维是指向各输入位置的概率，后$|C|$维是各类型标记的概率，$\kappa\in[0,1]$为超参数，混合上下文表示与原始词嵌入。这是第\ref{sec:nlp-analysis-paradigm-generation}小节“只含位置指针和类型标记”的彻底复制形式：输出不经过词表，因此不会生成原文没有的字。

训练使用教师强制，即每步以正确前缀为条件，最小化目标序列$\bm y^*$的负对数似然：
\begin{equation}
  L_{\mathrm{gen}}=
  -\sum_{u=1}^{|\bm y^{*}|}
  \log p_\theta(y_u^{*}\mid \bm y_{<u}^{*},x,c),
  \label{eq:nlp-extraction-generation}
\end{equation}
$c$为可选的任务条件，BARTNER中为空，下一小节的UIE中为模式提示。推理用束搜索逐步生成，再按规则解析：读取指针直到遇到类型标记，把指针两两配成片段。指针个数为奇数、片段起点大于终点或后一段起点不大于前一段终点时，该实体不合法。这些规则也可以写成解码时的掩码，例如类型标记之前必须有偶数个指针。

指针序列的位置信息可以直接还原；不足在于实体之间本无顺序，按原文位置排列是额外约定，模型生成正确集合的另一种排列也会受交叉熵惩罚，逐步生成也慢于网格或跨度方法的一次并行预测。

\subsubsection{UIE结构化生成与模式提示}
\label{sec:nlp-extraction-uie}

BARTNER只处理实体。若实体、关系和事件都要用同一个模型抽取，就需要一种能表达类型、片段及其关联的统一输出语言，以及一种告诉模型“这次要抽什么”的输入方式。Lu等的UIE为此设计了两部分\scite{lu2022uie}。

\term{结构化抽取语言}（structured extraction language, SEL）把抽取结果写成括号表达式。一条记录由若干“定位”组成，每个定位写出类型名与原文片段，下面可挂若干“关联”，每个关联写出角色名与片段：
\[
  \begin{aligned}
    \langle\text{记录}\rangle&::=\texttt{(}\ \{\langle\text{定位}\rangle\}\ \texttt{)},\\
    \langle\text{定位}\rangle&::=\texttt{(}\ \text{类型名}\ \texttt{:}\ \text{片段}\ \{\langle\text{关联}\rangle\}\ \texttt{)},\\
    \langle\text{关联}\rangle&::=\texttt{(}\ \text{角色名}\ \texttt{:}\ \text{片段}\ \texttt{)},
  \end{aligned}
\]
花括号表示可重复零次或多次。
实体只用定位，关系是带一个关联的定位，事件是触发词定位下挂多个论元关联。\term{结构模式提示}（structural schema instructor, SSI）把本次要识别的类型和关联名放在输入前面，用特殊标记\code{[spot]}引出类型名、\code{[asso]}引出关联名、\code{[text]}引出原文。对通知抽取地点与组织时，输入为“\code{[spot]} 地点 \code{[spot]} 组织 \code{[text]} 机器学习读书会原定于……地点为海棠楼201。”，目标为\code{((地点: 海棠楼201))}。模式提示换成关系或事件的名称，同一模型就转为关系或事件抽取。

UIE以T5为骨架，编码器读入$[c;x]$，其中$c$为SSI，解码器生成SEL。下游微调使用式~\eqref{eq:nlp-extraction-generation}的教师强制损失。为了让模型在没见过的模式上也会用这套语言，UIE先在三类数据上预训练：文本与结构配对的数据$\mathcal D_{\mathrm{pair}}$，来自知识库与百科文本的远程对齐；只有结构的数据$\mathcal D_{\mathrm{record}}$；以及无标注文本$\mathcal D_{\mathrm{text}}$。三项损失相加：
\begin{equation}
  \begin{aligned}
    L_{\mathrm{pre}}&=L_{\mathrm{pair}}+L_{\mathrm{record}}+L_{\mathrm{text}},\qquad
    L_{\mathrm{pair}}=-\sum_{(x,y)\in\mathcal D_{\mathrm{pair}}}\log p_\theta(y\mid x,c_{\mathrm{meta}}),\\
    L_{\mathrm{record}}&=-\sum_{y\in\mathcal D_{\mathrm{record}}}\sum_{u}\log p_\theta(y_u\mid y_{<u}),\qquad
    L_{\mathrm{text}}=-\sum_{x\in\mathcal D_{\mathrm{text}}}\log p_\theta(x''\mid x').
  \end{aligned}
  \label{eq:nlp-extraction-uie-pretrain}
\end{equation}
$L_{\mathrm{pair}}$教模型从文本到结构的映射，其提示$c_{\mathrm{meta}}$除了样本中真实出现的类型，还随机加入未出现的负类型，让模型学会对提示中的类型说“没有”；$L_{\mathrm{record}}$只训练解码器，把它当作结构语言的语言模型，学会生成合法的SEL；$L_{\mathrm{text}}$是T5式的跨度损坏去噪（第\ref{sec:nlp-llm-objective-t5-bart}小节），$x'$为遮去若干片段的文本，$x''$为被遮去的片段，用来保持通用语言能力。微调时UIE还在目标中随机插入\code{(类型: [NULL])}一类负定位，进一步训练拒绝输出。

模式提示只是条件，不能保证输出合法。推理时可以维护SEL的文法状态并施加约束：类型名位置只允许SSI中列出的名称，片段位置只允许原文的子串（用原文的前缀树检查），括号必须配对；机制见第\ref{sec:nlp-inference-grammar}小节。生成结束后还要把片段对齐回原文：若“海棠楼201”出现两次，字符串本身不能确定是哪一次，需要按上下文选择或保留全部候选位置。

与BARTNER的位置指针相比，UIE输出文字片段，换来跨任务的统一接口，代价正是这一步原文对齐。

\subsection{少样本与提示式实体识别}
\label{sec:nlp-extraction-ner-fewshot}

新类型往往只有几个标注样本，前述监督模型无法直接训练。度量学习路线用少量样本构造类别表示，对查询词元按距离分类，仍需在源领域上训练表示；大语言模型提示路线把类型定义和示例写进输入，不更新参数，但要处理格式解析、原文对齐与输出不稳定。

\subsubsection{原型网络与近邻标注}
\label{sec:nlp-extraction-prototype}

少样本实体识别通常按\emph{$N$类$K$样本}设定评价：每个测试任务给出$N$个新类型，每类只有约$K$个标注实体，组成\term{支持集}（support set）$\mathcal S$；模型要标注另一批句子组成的\term{查询集}（query set）。Few-NERD按句子采样任务，保证每类实体在支持集中出现$K$到$2K$次，并区分两种划分：训练与测试的粗类不同，或粗类相同而细类不同\scite{ding2021fewnerd}。词元级的标签空间是$N$个类型再加O。

\term{原型网络}（prototypical network）用每类支持样本的均值作为该类的代表\scite{ket-snell2017-prototypical}。设编码器把词元$x$映射为$f_\theta(x)\in\mathbb R^{d}$，$\mathcal S_t$为支持集中标签为$t$的全部词元，类$t$的原型和查询词元的类别分布为
\begin{equation}
  \bm c_t=\frac{1}{|\mathcal S_t|}\sum_{x'\in\mathcal S_t}f_\theta(x'),\qquad
  p(y=t\mid x)=\frac{\exp\bigl(-\lVert f_\theta(x)-\bm c_t\rVert_2^2\bigr)}{\sum_{t'}\exp\bigl(-\lVert f_\theta(x)-\bm c_{t'}\rVert_2^2\bigr)}.
  \label{eq:nlp-extraction-prototype}
\end{equation}
训练采用\term{元训练}（meta-training）：在源领域的已知类型上反复采样与测试同形的$N$类$K$样本任务，以查询词元的交叉熵更新$\theta$，使编码器学会把同类词元聚在一起；测试时冻结$\theta$，只用新类型的支持集计算原型。原型网络的一般讨论可参考\BookChapterRef{chap:ket-transfer-adaptation}{卷八的知识迁移章}。

用于实体识别时，O类暴露出原型的弱点。O包含标点、动词、普通名词等各种词元，表示空间里是多个分散的簇，均值可能落在一片空白区域。设二维教学表示中O的两个支持词元为$(0,1)$与$(0,-1)$，原型为$(0,0)$；LOC的支持词元为$(1,0)$，原型就是它本身。一个LOC查询词元位于$(0.45,0)$，到O原型的平方距离为$0.2025$，到LOC原型为$0.3025$，于是被判为O。若改看最近的单个支持词元，它到两个O词元的平方距离都是$1.2025$，到LOC词元为$0.3025$，判为LOC。

Yang与Katiyar据此提出NNShot：先在源领域上按普通监督方式训练一个实体识别模型，取其表示作为$f$，测试时把查询词元归到最近的支持词元所属的类，$\hat y(x)=\argmin_t\min_{x'\in\mathcal S_t}\lVert f(x)-f(x')\rVert_2^2$\scite{yang2020structshot}。它不需要元训练，近邻规则也不假设每类只有一个簇。逐词元分类仍可能输出不合法或不连贯的标签，StructShot再加入标签转移：新领域的类型与源领域不同，具体类型间的转移无法估计，于是把标签抽象为O、“与前一个同类型的I”和“与前一个不同类型的I”等少数几类，在源领域上估计这些\emph{抽象转移}的概率，再展开到新类型上。解码用Viterbi求
\begin{equation}
  \hat{\bm y}=\argmax_{\bm y}\sum_{i=1}^{n}\Bigl[\log p_{\mathrm{NN}}(y_i\mid x_i)+\lambda\log p_{\mathrm{abs}}(y_i\mid y_{i-1})\Bigr],
  \label{eq:nlp-extraction-structshot}
\end{equation}
$p_{\mathrm{NN}}$由各类最近支持词元的距离经softmax得到。原文把发射概率与转移概率相乘，由于两者分别估计、尺度不一致，再用温度超参数重新归一化转移概率；式中的$\lambda$是本书对这一尺度调节的对数形式写法，不是原文的参数。它与第\ref{sec:nlp-extraction-bio}小节的CRF形式相同，区别在于发射来自近邻距离、转移来自源领域的抽象统计，两者都不在新类型上训练。

三者的取舍在于怎样用几个样本定义类别：原型网络适合类内紧凑的情形；NNShot适应多峰的O类，但对支持集中的个别错误标注更敏感；StructShot补回标签依赖，前提是源领域与目标领域的标签转移模式相近。

\subsubsection{大语言模型提示式抽取}
\label{sec:nlp-extraction-llm-ner}

大语言模型不需要元训练，也不需要支持集上的编码器。上下文学习（第\ref{sec:nlp-inference-icl-shots}小节）把任务说明与少量示例写进提示，模型在不更新参数的情况下给出输出。提示式实体识别的提示$P$由三部分组成：指令$\mathcal I$写明类型定义、边界规则和“没有就返回空”等要求；$k$个示例$(x^{(1)},y^{(1)}),\ldots,(x^{(k)},y^{(k)})$展示输入与输出的对应；最后是待标注文本$x$与输出格式说明。模型按$\hat y\sim p_{\mathrm{LM}}(\cdot\mid P(\mathcal I,\{(x^{(r)},y^{(r)})\},x))$生成，再由解析程序把$\hat y$转成实体集合。

示例与输出格式是两个主要设计点。示例数受上下文长度限制，常按与$x$的表示相似度检索最近的$k$个训练句，让示例覆盖相近的边界与类型模式。输出格式有两类。一类是结构化列表，例如JSON数组\code{[\{"text": "海棠楼201", "type": "地点"\}]}，便于解析，可用JSON Schema约束解码保证格式合法（第\ref{sec:nlp-inference-grammar}小节）；但它只给字符串，需要在原文中查找位置，同一字符串出现多次时还要消歧。另一类是标记改写，让模型把原文抄写一遍并在实体两侧加记号，例如GPT-NER用\code{@@海棠楼201\#\#}标出地点，每次只问一种类型\scite{wang2023gptner}；输出与原文逐字对齐即可恢复偏移，代价是要抄写全文，长文本上成本高，抄写偏差也会导致对齐失败。

单次生成的结果不稳定：同一提示多次采样，可能漏掉或多出某个实体，边界也可能漂移。可以借用自一致性（第\ref{sec:nlp-execution-self-consistency}小节）的思路做实体级投票。采样$R$次得到实体集合$\hat{\mathcal E}^{(1)},\ldots,\hat{\mathcal E}^{(R)}$，对每个候选“片段加类型”$e$计票，保留票数达到阈值$\kappa$的实体：
\begin{equation}
  \operatorname{vote}(e)=\sum_{r=1}^{R}\mathbb 1\bigl[e\in\hat{\mathcal E}^{(r)}\bigr],\qquad
  \hat{\mathcal E}=\{e:\operatorname{vote}(e)\geq\kappa\}.
  \label{eq:nlp-extraction-entity-vote}
\end{equation}
与第\ref{sec:nlp-execution-self-consistency}小节对单个答案投票不同，这里每个实体单独投票，输出仍是集合。教学示例：$R=5$，“海棠楼201／地点”出现4次，边界不全的“海棠楼／地点”出现1次，“机器学习读书会／组织”出现2次；取$\kappa=3$时只保留第一项。读书会是否算组织取决于指令中的类型定义，票数分散正反映了模型对定义的不确定，补充定义比提高采样次数更直接。另一种做法是GPT-NER采用的自我核验：对每个候选实体再问一次“原文中的‘海棠楼201’是不是地点”，只保留肯定回答，用来抑制模型把非实体标成实体的倾向。无论哪种做法，输出不合格式、片段不是原文子串、类型不在指令列出的集合中，都应计为错误，不能静默修补；原文中没有的“李老师”即使格式合法，也是无依据的新增。

提示式抽取几乎不需要标注，换类型只需改指令，适合冷启动与类型频繁变化的场景；代价是长提示与多次采样的成本，严格边界上的准确率通常也低于在充足标注上训练的专用模型。两者可以衔接：用提示式抽取产生初始标注，经人工校对后训练第\ref{sec:nlp-extraction-ner-span}小节的跨度模型；或者把多种抽取任务整理为指令数据，对模型做指令微调（第\ref{chap:nlp-supervised-finetuning}章），使它在抽取格式与边界上更稳定。

\section{实体链接}
\label{sec:nlp-extraction-linking}

识别出提及以后，还要确定它对应知识库中的哪个对象。\term{实体链接}（entity linking, EL）把提及映射到给定版本知识库中的条目；没有可确认的对应项时输出NIL。为了比较本节的四种方法，构造两个同名“海棠大学”：$e_1$是教学大学，描述含“计算机学院”；$e_2$是教学小说中的虚构学校，描述含“故事人物”。“海棠大学计算机学院发布通知”应链接$e_1$，“小说中的海棠大学迎来新角色”应链接$e_2$。

实体链接有四项挑战。别名多：同一对象有全称、简称和俗称；同名多：一个名字对应多个对象；知识库大：候选对象可达数百万，不能对每个都做精细比较；库外实体：提及的对象可能尚未登记。表~\ref{tab:nlp-extraction-linking-methods}按这些挑战对照本节的四种方法：别名字典与先验、BLINK召回与重排、文档全局一致性，以及GENRE生成式名称预测。

\begin{table*}[htbp]
  \centering
  \small
  \begin{tabularx}{\linewidth}{@{}>{\raggedright\arraybackslash}p{22mm}>{\raggedright\arraybackslash}p{30mm}>{\raggedright\arraybackslash}X>{\raggedright\arraybackslash}p{30mm}>{\raggedright\arraybackslash}p{34mm}@{}}
    \toprule
    方法 & 候选来源 & 使用的证据 & 新增实体 & 主要代价\\
    \midrule
    别名字典与先验 & 登记名称与锚文本计数 & 链接先验、上下文相似、类型兼容 & 需补别名与计数 & 漏别名则无法召回\\
    BLINK & 实体向量近邻检索 & 提及上下文与实体描述的语义匹配 & 写好描述即可编码入库 & 索引维护；重排需逐候选联合编码\\
    全局一致性 & 前一阶段的候选 & 局部分数加实体间兼容 & 依赖实体表示 & 联合推断需近似；跨领域文档易误伤\\
    GENRE & 名称前缀树 & 上下文条件下的名称生成概率 & 加入名称即合法 & 名称须唯一；串行解码\\
    \bottomrule
  \end{tabularx}
  \caption{实体链接的方法谱系。前两种方法常组成“召回—重排”流水线，全局一致性作用在重排之后，GENRE以一个生成模型同时完成召回与选择。}
  \label{tab:nlp-extraction-linking-methods}
\end{table*}

评价要分两段看。候选召回问正确身份是否进入候选集合$C(m)$，常用前$k$召回率；重排问在集合中是否选对。正确身份一旦在召回阶段遗漏，后续再准确的重排也无法恢复。评价还须声明知识库快照、提及是否给定、候选规模和NIL比例。

\subsection{别名字典与实体先验排序}
\label{sec:nlp-extraction-alias}

最直接的候选来源是登记名称、人工别名和历史链接记录。百科文本中的锚文本是天然的链接标注：若“海棠大学”这段锚文本有$c(m,e)$次指向条目$e$，则\term{链接先验}（link prior）为$p(e\mid m)=c(m,e)/\sum_{e'}c(m,e')$，它表示不看语境时这个名字指向$e$的频率。别名表把“海棠大学”映射到$\{e_1,e_2\}$，再加入上下文与类型信息组成线性评分：
\begin{equation}
  s(m,e)=\log p(e\mid m)+
  \lambda\,\operatorname{sim}(x_m,d_e)+
  \eta\,\operatorname{type}(m,e),
  \label{eq:nlp-extraction-alias}
\end{equation}
其中$x_m$为提及上下文，$d_e$为实体描述，$\operatorname{sim}$为二者的相似度（如词袋余弦），$\operatorname{type}$为提及类型与实体类型的兼容分数，$\lambda,\eta$在开发集上选择。先验只能从训练期的计数估计，罕见别名需要平滑。

若教学先验为$(0.8,0.2)$，小说语境的相似度为$(0,2)$，取$\lambda=1$且类型项相等，两候选的分数为$\log0.8\approx-0.2231$与$\log0.2+2\approx0.3906$，选$e_2$，语境推翻了先验的首选；在学院语境中相似度换成$(2,0)$，则选$e_1$。这个方法查表即可、速度很快，先验本身也是强基线；但候选质量完全取决于别名覆盖，若别名表漏掉$e_2$，任何语境都选不到它。低分拒绝为NIL的阈值需要在含库外提及的开发集上校准。

\subsection{BLINK双编码召回与交叉编码重排}
\label{sec:nlp-extraction-blink}

别名表漏候选、只用浅层语境，两项不足都来自“按名字找”。Wu等的BLINK改为“按意思找”：用语义向量在全库中检索候选，再用更精细的模型重排\scite{wu2020blink}。它的设计目标是\emph{零样本链接}：测试时的实体在训练中从未出现，只有一段文字描述。

第一阶段是\term{双编码器}（bi-encoder）。提及一侧的输入是带边界标记的上下文“$[\mathrm{CLS}]$ 左文 $[\mathrm{M_s}]$ 提及 $[\mathrm{M_e}]$ 右文 $[\mathrm{SEP}]$”，实体一侧的输入是“$[\mathrm{CLS}]$ 名称 $[\mathrm{ENT}]$ 描述 $[\mathrm{SEP}]$”，两个独立的BERT编码器各取$[\mathrm{CLS}]$位置的输出，得到$\bm u_m,\bm v_e\in\mathbb R^{d}$，相似分数为内积$s_{\mathrm{bi}}(m,e)=\bm u_m^\top\bm v_e$。训练使用第\ref{chap:nlp-classification-matching-organization}章式~\eqref{eq:nlp-matching-infonce}的对比损失，温度取1，对正确实体$e^+$与负实体集合$N_m$，
\begin{equation}
  L_{\mathrm{bi}}=
  -\log\frac{\exp s_{\mathrm{bi}}(m,e^+)}
  {\sum_{e\in\{e^+\}\cup N_m}\exp s_{\mathrm{bi}}(m,e)}.
  \label{eq:nlp-extraction-blink}
\end{equation}
负例先取同一批中其他提及的正确实体，这样一批$B$个提及只需编码$B$个实体，就得到$B\times B$的分数矩阵；第二轮再用训练好的双编码器检索出分数高的错误实体作为\term{难负例}（hard negative）继续训练。两侧编码相互独立，因此全部实体向量可以离线算好、建近邻索引，查询时只编码提及一次。

第二阶段是\term{交叉编码器}（cross-encoder）。它把提及上下文与候选实体的名称和描述拼成一个输入，送入同一个BERT，取$[\mathrm{CLS}]$输出$\bm g_{m,e}\in\mathbb R^{d}$，打分$s_{\mathrm{cross}}(m,e)=\bm w^{\top}\bm g_{m,e}$。训练时在双编码器召回的前$k$个候选上做softmax交叉熵。联合编码让提及中的每个词元都能注意到描述中的每个词元，比较更细致，但每个候选都要重新编码一遍，只能用于少量候选。

学院语境的教学计算：令$\bm u_m=(1,0.5)^\top$，$\bm v_{e_1}=(0.8,0.4)^\top$，$\bm v_{e_2}=(0.9,0.5)^\top$，内积为1.00与1.15，双编码器把错误的$e_2$排在前面。若只召回一项，$e_1$已经丢失；召回两项后，设交叉编码分数为$(2.4,0.2)$，则$p(e_1\mid C(m),x)=e^{2.4}/(e^{2.4}+e^{0.2})\approx0.9002$，重排纠正了召回的次序。这里的概率只在候选集合内归一化，不代表“库中没有对应项”的概率，NIL需要另设阈值或拒绝分类器。

两段的代价不同。精确检索每个提及为$O(|\mathcal K|d)$，$|\mathcal K|$为库中实体数，近似近邻索引以少量召回损失换取查询速度；重排对$k$个候选各做一次完整编码。新实体只需写好描述、编码后加入索引，这是它能做零样本链接的原因；编码器更新后则要重算全部实体向量。

\subsection{文档全局一致性重排}
\label{sec:nlp-extraction-global-linking}

逐个提及独立链接，忽略了同一文档中的实体往往彼此相关。“海棠大学”与“计算机学院”出现在同一通知中，若后者明确链接到校园的学院，前者也更可能是校园的大学。把这种相关性写进目标，就得到全局链接。设文档有提及$m_1,\ldots,m_q$，联合赋值为$\bm e=(e_1,\ldots,e_q)$，$\Psi_i(e_i)$为第$i$个提及的局部分数，$\Phi(e_i,e_j)$为两实体的兼容分数：
\begin{equation}
  \widehat{\bm e}=\argmax_{\bm e\in\prod_i C(m_i)}
  \left[\sum_{i}\Psi_i(e_i)+
  \frac{2}{q-1}\sum_{i<j}\Phi(e_i,e_j)\right].
  \label{eq:nlp-extraction-global-linking}
\end{equation}
这是一个以提及为节点、候选实体为取值的全连接条件随机场，一元势为局部分数，成对势为兼容分数；系数$2/(q-1)$使成对项的总量不随提及数增长。

Ganea与Hofmann给出了一种端到端的实现\scite{ganea2017deep}。局部分数用实体嵌入$\bm x_e$对上下文词做注意力：每个上下文词$w$先按它与候选实体的最大相关度打分，只保留得分最高的若干词，归一化为权重$\beta(w)$，局部分数为$\Psi(e)=\bm x_e^\top\bm B\sum_w\beta(w)\bm x_w$，$\bm B$为对角矩阵；再用一个小型网络把它与$\log p(e\mid m)$组合。成对势为$\Phi(e,e')=\bm x_e^\top\bm C\bm x_{e'}$，$\bm C$也是对角矩阵。全连接图上精确求解式~\eqref{eq:nlp-extraction-global-linking}的代价为$O(k^q)$，$k$为每个提及的候选数，因此改用截断的\term{环置信传播}（loopy belief propagation）：在有环图上照常迭代传递消息，固定轮数后停止，以此近似每个提及的边缘分布。可参考\BookChapterRef{chap:pgm-variational}{卷六的变分推断与确定性近似章}，那里说明了消息传递在有环图上为何只是近似、何时可能不收敛。第$t$轮从提及$i$传给提及$j$的最大积消息为
\begin{equation}
  \mu^{(t+1)}_{i\to j}(e)=\max_{e'\in C(m_i)}\Bigl[\Psi_i(e')+\tfrac{2}{q-1}\Phi(e',e)+\sum_{k\neq i,j}\mu^{(t)}_{k\to i}(e')\Bigr],
  \label{eq:nlp-extraction-lbp}
\end{equation}
它表示“在$m_j$取$e$的条件下，$m_i$及其上游能贡献的最大分数”。Ganea与Hofmann还对消息做阻尼：把本轮消息归一化后与上一轮消息按系数$\delta\in(0,1]$加权平均，以稳定迭代。$T$轮后，提及$i$的置信为$b_i(e)=\Psi_i(e)+\sum_{k\neq i}\mu^{(T)}_{k\to i}(e)$，经softmax后再与先验组合成最终分数$\rho_i(e)$。整个展开的消息传递过程可微，训练以排序间隔损失$\sum_i\sum_{e\neq e_i^*}\max(0,\gamma-\rho_i(e_i^*)+\rho_i(e))$端到端更新嵌入、注意力和组合网络。

教学数值：大学的局部分数为校园0.6、小说0.8，学院为校园0.9、小说0.4；成对项（已含系数$2/(q-1)$）同领域取0.7，跨领域取0。四种组合的总分为校园—校园2.2、校园—小说1.0、小说—校园1.7、小说—小说1.9，全局选择校园组合，改变了大学的局部首选。用式~\eqref{eq:nlp-extraction-lbp}计算，学院传给大学的消息为：大学取校园时$\max(0.9+0.7,\ 0.4+0)=1.6$，取小说时$\max(0.9+0,\ 0.4+0.7)=1.1$；大学的置信为校园$0.6+1.6=2.2$、小说$0.8+1.1=1.9$，选校园。两个节点的图没有环，一轮消息即得到与穷举相同的答案；提及增多后图中有环，消息传递就只是近似。

Le与Titov进一步认为实体之间的兼容不止一种：两个实体可能同属一个机构，也可能是人物与其所在城市。他们引入$K$种隐关系，成对势写成$\Phi(e_i,e_j)=\sum_{r=1}^{K}a_{i,j,r}\,\bm x_{e_i}^\top\bm R_r\bm x_{e_j}$，每种关系有自己的对角矩阵$\bm R_r$，权重$a_{i,j,r}$由两个提及的上下文表示计算并归一化，无需关系标注\scite{le2018latent}。

全局一致性依赖“同一文档的实体相互关联”这一假设。一篇专门比较现实大学与小说学校的文章本来就跨领域，强制一致会把正确的跨领域链接改错；成对势的权重需要在包含这类文档的验证集上调整，NIL候选也应参与联合推断。

\subsection{生成式实体名称预测}
\label{sec:nlp-extraction-generative-linking}

BLINK把实体压成向量，需要为数百万实体存储并维护索引，还需要负例采样训练。De Cao等的GENRE换一种思路：每个实体有唯一的规范名称（如百科标题），链接就是在上下文条件下生成这个名称\scite{decao2021genre}。模型是BART，输入是用$[\mathrm{START\_ENT}]$与$[\mathrm{END\_ENT}]$标出提及的上下文$x_m$，输出名称的词元序列$\bm y_e$，评分为$\log p_\theta(\bm y_e\mid x_m)=\sum_u\log p_\theta(y_{e,u}\mid\bm y_{e,<u},x_m)$，用已链接提及按式~\eqref{eq:nlp-extraction-generation}训练。生成的名称必须是库中存在的名称，GENRE为此把全部名称建成一棵前缀树，在束搜索（第\ref{sec:nlp-generation-beam}小节）的每一步只允许能延续成某个合法名称的词元，机制与第\ref{sec:nlp-inference-grammar}小节的约束解码相同。

教学知识库把两个对象规范化为“海棠大学（校园）”和“海棠大学（小说）”。生成“海棠大学（”之后，前缀树只允许“校园）”与“小说）”两条延伸；若两条完成路径的概率分别为0.12与0.08，就选前者，再查名称表恢复$e_1$。输入中的别名可以相同，输出名称必须一一对应身份，知识库中的重名标题需要消歧后缀。

与双编码器相比，GENRE的存储只是一棵前缀树，不需要负例采样；生成时每个词元都能与提及上下文充分交互，却不必像交叉编码器那样逐个候选编码。它的局限是串行解码，束宽决定了能比较的候选数；新增实体只要加入前缀树就成为合法输出，但模型未必学会在新名称之间作出选择。无对应项时需要另设NIL输出或拒绝判断。
